%\documentclass[12pt,letterpaper]{article}



%%%
%% Required packages:
%%
\usepackage{etex}
\usepackage{amsmath}
\usepackage{amssymb}
\usepackage{amstext}
\usepackage{amsthm}
\usepackage{fancyhdr,fancybox}
\usepackage{mathrsfs} % need for \mathscr command
\usepackage{multicol}
\usepackage[normalem]{ulem}
%\usepackage{pictex,pictexplus}
% \usepackage{pictexwd}
\usepackage{rotating}
\usepackage{url}
\usepackage{xfrac}
\usepackage{booktabs}
\usepackage{color,soul}
\usepackage{comment}

\usepackage{hyperref}
\hypersetup{
    colorlinks=true,     % false: boxed links; true: colored links
    linkcolor=blue,      % color of internal links
    citecolor=blue,      % color of links to bibliography
    filecolor=blue,      % color of file links
    urlcolor=blue        % color of external links
}


\usepackage{tikz}
\usetikzlibrary{backgrounds,calc}

%%
%% Common formatting commands
%%
\setlength{\parindent}{0in}
\setlength{\textwidth}{7in}
\setlength{\evensidemargin}{-0.25in}
\setlength{\oddsidemargin}{-0.25in}
\setlength{\parskip}{.5\baselineskip}
\setlength{\topmargin}{-0.5in}
\setlength{\textheight}{9in}

%               Problem and Part
\newcounter{problemnumber}
\newcounter{partnumber}
\newcounter{subpartnumber}
\newcommand{\Problem}{%
  \refstepcounter{problemnumber}%
  \setcounter{partnumber}{0}%
  \setcounter{subpartnumber}{0}%
  \item[\makebox{\hfil\textbf{\theproblemnumber.}\hfil}]%
  }
\newcommand{\InProblem}[2][1.6]{%
  \refstepcounter{problemnumber}%
  \setcounter{partnumber}{0}%
  \setcounter{subpartnumber}{0}%
  \textbf{\theproblemnumber.}\ \ \parbox[t]{#1in}{#2}%
}
\newcommand{\InSmallProblem}[1]{\InProblem[1.05]{#1}}
\newcommand{\NoProblem}[1]{%
  \mbox{\hfil\textbf{#1}\hfil}%
  }
\newcommand{\UnProblem}{%
  \refstepcounter{problemnumber}%
  \setcounter{partnumber}{0}%
  \setcounter{subpartnumber}{0}%
  \mbox{\hfil\textbf{\theproblemnumber.}\hfil}%
  }
\newcommand{\InFlexProblem}[1]{%
  \refstepcounter{problemnumber}%
  \setcounter{partnumber}{0}%
  \setcounter{subpartnumber}{0}%
  \textbf{\theproblemnumber.}\ \ #1%
}
%%  Part:
\newcommand{\Part}{%
  \renewcommand{\thepartnumber}{(\alph{partnumber})}%
  \refstepcounter{partnumber}
  \setcounter{subpartnumber}{0}%
  \item[(\alph{partnumber})]%
}
\newcommand{\InPart}[2][1.75]{%
  \renewcommand{\thepartnumber}{(\alph{partnumber})}%
  \refstepcounter{partnumber}%
  \setcounter{subpartnumber}{0}%
  (\alph{partnumber})\ \ \parbox[t]{#1in}{#2}%
}
\newcommand{\UnPart}{%
  \refstepcounter{partnumber}%
  \renewcommand{\thepartnumber}{(\alph{partnumber})}%
  \setcounter{subpartnumber}{0}%
  (\alph{partnumber})%
}
\newcommand{\InLargePart}[1]{\InPart[3.05]{#1}}
\newcommand{\InFullPart}[1]{\InPart[6.2]{#1}}
\newcommand{\InSmallPart}[1]{\InPart[1.05]{#1}}
%% SubPart:
\newcommand{\SubPart}{%
  \refstepcounter{subpartnumber}%
  \item[(\roman{subpartnumber})]%
  }
\newcommand{\InSubPart}[2][2.25]{%
  \stepcounter{subpartnumber}%
  (\roman{subpartnumber})\ \ \parbox[t]{#1in}{#2}%
}
\newcommand{\InSmallSubPart}[1]{\InSubPart[1.5]{#1}}
\newcommand{\InTinySubPart}[1]{%
  \stepcounter{subpartnumber}%
  (\roman{subpartnumber})\ \ \makebox{#1}%
}
\newcommand{\InMySubPart}[2][2.25]{%
  \stepcounter{subpartnumber}%
  \parbox[t]{#1in}{\makebox[0.3in][l]{(\roman{subpartnumber})}\ \ {#2}}%
}
\newcommand{\TrueFalse}{\textbf{T}\ \ \textbf{F}\ \ }
\newcommand{\TFPart}[1]{% 
  \stepcounter{partnumber}%
  \setcounter{subpartnumber}{0}%
  \item[(\alph{partnumber})] \TrueFalse\parbox[t]{5.5in}{#1}}

\newcommand{\Points}[1]{(#1 points) \ }
\newcommand{\Point}[1]{(#1 point) \ }


%% For Answers:
\newcounter{answernumber}
\newcommand{\Answer}{%
  \refstepcounter{answernumber}%
  \setcounter{partnumber}{0}%
  \setcounter{subpartnumber}{0}%
  \item[\makebox{\hfil\textbf{\theanswernumber.}\hfil}]%
  }


% Setting up the headers for the first page
%\chead{} % will default to what is typed in the file.
\fancyhead[L]{\textsc{Math 22a}}
\fancyhead[R]{\textsc{Fall 2024}}
\fancyfoot[R]{
	\vspace*{\fill}\footnotesize\textcopyright{} \emph{Harvard Math 22a}	
}
\renewcommand{\headrulewidth}{0pt}

%\fancyhead[C]{}
%	\chead{}	
%	\lhead{}
%	\rhead{}
%	\cfoot{}


% Putting copyright on the footer of each page
%\fancypagestyle{secondpage}{
%	\fancyhead[C]{TESTing again} % change this to be blank
%		%\chead{}	
%	\fancyhead[L]{bears} % change this to be blank
%	\fancyhead[R]{dogs} % change this to be blank
%	\fancyfoot[R]{ %keeps the footer the same
%		\vspace*{\fill}\footnotesize\textcopyright{} \emph{Harvard Math 22a}	
%	}
%}
%\renewcommand{\headrulewidth}{0pt}
%\pagestyle{fancy}
\pagestyle{empty}


%\lhead{}
%\rhead{}
%\rfoot{\vspace*{\fill}\footnotesize\textcopyright{} \emph{Harvard Math 22a}}
%\renewcommand{\headrulewidth}{0pt}
%\pagestyle{fancy}




%%
%% Common shortcut commands:
%%
\newcommand{\ds}{\displaystyle}
\newcommand{\sgn}{\operatorname{sgn}}
\renewcommand{\Pr}{\operatorname{P}}
\newcommand{\CPr}[3][{}]{\Pr_{\text{#1}}\left(#2\,|\,#3\right)}

\newcommand{\C}{\mathbb{C}}
\newcommand{\R}{\mathbb{R}}
\newcommand{\Z}{\mathbb{Z}}
\newcommand{\N}{\mathbb{N}}
\newcommand{\Q}{\mathbb{Q}}
\newcommand{\D}{\mathbb{D}}

\newcommand{\rref}{\operatorname{rref}}
\newcommand{\im}{\operatorname{im}}
\newcommand{\rank}{\operatorname{rank}}
\newcommand{\diag}{\operatorname{diag}}
\newcommand{\Span}{\operatorname{span}}
%\renewcommand{\vec}[1]{\mathbf{#1}}
\newcommand{\charpoly}{\mathscr{P}}
\newcommand{\nicefrac}[2]{\sfrac{#1}{#2}} % using xfrac package
\newcommand{\topology}{\mathcal{T}}
\newcommand{\Inn}{\operatorname{Inn}}

\newcommand{\blankbox}[2]{\fbox{\rule{#1}{0in}\rule{0in}{#2}}}

\newenvironment{myindentpar}[1]%
 {\begin{list}{}%
         {\setlength{\leftmargin}{#1}
          \setlength{\rightmargin}{\leftmargin}}%
         \item[]%
 }
 {\end{list}}
\newtheorem{theorem}{Theorem}
\newtheorem*{NStheorem}{Theorem}
\newtheorem{cor}[theorem]{Corollary}
\newtheorem*{claim}{Claim}
\newtheorem{defn}{Definition}

\newenvironment{amatrix}[1]{%
  \left(\begin{array}{@{}*{#1}{c}|c@{}}
}{%
  \end{array}\right)
}

%--------grstep
% For denoting a Gauss' reduction step.
% Use as: \grstep{\rho_1+\rho_3} or \grstep[2\rho_5 \\ 3\rho_6]{\rho_1+\rho_3}
\newcommand{\grstep}[2][\relax]{%
   \ensuremath{\mathrel{
       {\mathop{\longrightarrow}\limits^{#2\mathstrut}_{
                                     \begin{subarray}{l} #1 \end{subarray}}}}}}
\newcommand{\swap}{\leftrightarrow}


%\usetikzlibrary{arrows}
%\usepackage{wrapfig}

%\makeatletter
%\renewcommand*\env@matrix[1][*\c@MaxMatrixCols c]{%
%	\hskip -\arraycolsep
%	\let\@ifnextchar\new@ifnextchar
%	\array{#1}}
%\makeatother
%
%\def\env@matrix{\hskip -\arraycolsep
%	\let\@ifnextchar\new@ifnextchar
%	\array{*\c@MaxMatrixCols c}}

\section{{Linear Independence, $\vec\S$4.3}}% 
\chead{\Large\textbf{Linear Independence, $\vec\S$4.3}}% 

%\begin{document}
\thispagestyle{fancy}

Recall from last class:

\begin{itemize}
\item A real {\bf vector space} is a non-empty set $V$ on which two operations, scalar multiplication and addition, are defined subject to the following 10 axioms. We call elements in $V$ vectors. Let $u,v,w \in V$ and $c,d\in \mathbb{R}$.
\begin{enumerate}
\item $u+v \in V$.
\item $u+v=v+u$.
\item $(u+v)+w=u+(v+w)$.
\item there exists a vector $0$ in $V$ such that $u+0=u$.
\item there exists a vector $(-u)$ in $V$ such that $u+(-u)=0$.
\item $cu \in V$.
\item $c(u+v)=cu+cv$.
\item $(c+d)u=cu+du$.
\item $c(du)=(cd)u$.
\item $1\cdot u =u$.
\end{enumerate}
\item A {\bf subsapce} of a vector space $V$ is a subset of $V$ such that
\begin{enumerate}
\item $0_V \in H$.
\item $H$ is closed under vector addition; i.e. if $u, v \in H$, then $u+v \in H$.
\item $H$ is closed under scalar multiplication; i.e. if $c \in \mathbb{R}$ and $u \in H$, then $cu \in H$.
\end{enumerate}
\item {\bf Theorem 4.1} If $v_1, \dots, v_p \in V$, then $\text{Span}\{v_1 ,\dots, v_p\}$ is a subspace of $V$. 
\item We call $\text{Span}\{v_1 ,\dots, v_p\}$ the subspace generated by or spanned by $v_1,\dots, v_p$.
\item 	The {\bf null space} of an $m \times n$ matrix $A$ is $$\text{Null}(A)=N(A)=\{x \in \mathbb{R}^n : A\vec{x}=\vec{0}\}.$$
\item {\bf Theorem:} If $A$ is an $m\times n$ matrix, then $N(A)$ is a subspace of $\mathbb{R}^n$.
\item 	The {\bf column space} of an $m\times n$ matrix $A$ is the span of the columns of $A$. We denote the column space of $A$ by $\text{Col}(A)$.\
\item Let $V$ and $W$ be vector spaces, then $T: V \to W$ is a {\bf linear transformation} if 
\begin{enumerate}
	\item $T(u+v)=T(u)+T(v)$ for all $u, v \in V$ and 
	\item $T(c u)=c T(u)$ for all $c \in \mathbb{R}$ and $u \in V$.
\end{enumerate}
\item 	The {\bf range or image } of $T$ is defined by $$\text{Range}(T)=\text{im}(T)=\{T(x): x \in V\}=\{w\in W: \text{there exists } a \in V \text{ such that } T(a)=w\}.$$ 

The {\bf kernel} of $T$ is defined by $\text{ker}(T)=\{ x \in V: T(x)=0_W\}.$
\end{itemize}

\newpage

\begin{defn}
	The subset $\{ v_1, v_2, \dots, v_p\}$ of the vector space $V$ is {\bf linearly independent} if $$c_1 v_1+\dots+ c_p v_p =0_V$$ has only the trivial solution.
\end{defn}

{\bf Theorem:} Given $\{ v_1, \dots, v_p\}$ with $p\geq 2$ and $v_1\neq 0_V$. The vectors $v_1, \dots, v_p$ are linearly dependent if and only if for some $j>1$, the vector $v_j$ is a linear combination of $v_1, \dots, v_{j-1}$.

\begin{enumerate}


\Problem 
\Part Let $C[0,2\pi]$ be the vector space of continuous real-valued functions on the interval $[0,2\pi]$. Is $\{ \sin{t}, \cos{t} \}$ linearly independent?

\vfill

\Part Consider $1+t^2 \in \mathbb{P}_2$ and $1-t^2 \in \mathbb{P}_2$. Are they linearly independent?

\vfill

\newpage

\begin{defn} Let $H$ be a subspace of a vector space $V$. The set $\{v_1, \dots, v_p\}$ is a {\bf basis} for $H$ if
	\begin{enumerate}
		\item  $\{v_1, \dots, v_p\}$ is linearly independent,
		\item $H=\text{Span}\{v_1, \dots, v_p \}$.
	\end{enumerate}
\end{defn}

{\bf Spanning Set Theorem:} Let $S=\{ v_1,\dots,v_p\} \subseteq V$ and $H=\text{Span}\{v_1,\dots, v_p \}$.
\begin{enumerate}
	\item Suppose $v_k$ can be written as a linear combination of the other vectors in $S$, then $S-\{v_k\}$ still spans $H$.
	\item If $H\neq \{ 0_V\}$, some subset of $S$ is a basis for $H$.
\end{enumerate} 

\Problem Prove the theorem.

\newpage

{\bf Theorem:} The pivot columns of an $m \times n$ matrix $A$ form a basis for $\text{Col}(A)$.

\Problem Prove the theorem.

\vfill

\Problem Let $B= \begin{bmatrix} 1 & 0 & 6 & 5\\ 0 & 2 & 5 & 3\\ 0 & 0 &0 & 0\\ \end{bmatrix}$. Find a basis for $\text{Col}(B)$.

\vfill

\Problem Is $\{1+t, 1-t,2\}$ a basis for $\mathbb{P}_1$?

\vfill

%\newpage

%{\bf Replacement Theorem:} Let $V$ be a vector space that is generated by set $G$. Suppose $|G|=n$. Let $L$ be a linearly independent subset of $V$ with $|L|=m$. Then $m \leq n$ and there exists a subset $H$ of $G$ containing exactly $n-m$ vectors such that $H \cup L$ generates $V$.

%\Problem Prove the theorem.
%
%\vfill

%{\bf Corollary:} Let $V$ be a vector space with a finite basis, then every basis has the same number of vectors.

\end{enumerate}

\begin{comment}
\newpage
\chead{\Large\textbf{Linear Independence -- Answers / Solutions}}%
\lhead{}
\rhead{}
\thispagestyle{fancy}

\begin{enumerate}
  \Answer 
  \begin{enumerate}
  \item Consider the vector equation $$c_1 \sin{t}+c_2 \cos{t} = 0$$ for scalars $c_1, c_2 \in \mathbb{R}$. Note two functions are equal means they are equal for all values of $t$. Thus the equation must hold for particular $t$ values. If $t=0$, then we get the equation $c_1(0)+c_2 (1) =0$, and hence $c_2=0$. If $t=\pi/2$, then we get $c_1(1)+c_2(0)=0$, and hence $c_1=0$. Thus the only solution to $c_1 \sin{t}+c_2 \cos{t} = 0$ is $c_1=c_2=0$. Thus $\sin{t}$ and $\cos{t}$ are linearly independent.
  \item Consider the vector equation $$c_1 (1+t^2)+c_2 (1-t^2)=0,$$ and regrouping terms we get $$(c_1+c_2)+t^2(c_1-c_2)=0.$$ Now equating coefficients we get $c_1+c_2=0$ and $c_1-c_2=0$, and hence $c_1=c_2=0$. Thus $1+t^2$ and $1-t^2$ are linearly independent.
  \end{enumerate}
  \Answer
  
  \begin{proof}
  	Let $V$ be a vector space and let $S=\{v_1, \dots, v_p\}\subseteq V$. Let $H=\text{Span}\{v_1,\dots, v_p \}$. 
  	
  	{\bf (a)} Without loss of generality, we assume $v_p$ is a linear combination of $v_1,\dots, v_{p-1}$. Thus there exists $c_1, \dots, c_{p-1} \in \mathbb{R}$ such that $$v_p=c_1 v_1+\dots +c_{p-1} v_{p-1}.$$ Let $x \in \text{Span}\{v_1,\dots, v_p\}$, then there exists $d_1, \dots, d_p \in \mathbb{R}$ such that $$x=d_1v_1 +\dots+ d_p v_p.$$ Plugging in our expression for $v_p$ and using the vector space axioms we can simplify $x$ in the following way:
  	$$x =d_1v_1 +\dots+ d_p v_p =d_1v_1 +\dots+ d_p (c_1 v_1+\dots +c_{p-1} v_{p-1})=(d_1 +d_p c_1)v_1+\dots+ (d_{p-1}+d_pc_{p-1})v_{p-1}.$$
  	
  	Thus we can express $x$ as a linear combination of $v_1, \dots, v_{p-1}$, and hence $x \in \text{Span}\{v_1,\dots, v_{p-1}\}.$
  	
  	{\bf (b)} If $S$ is linearly independent, then $S$ is a basis for $H$, and hence the claim holds.
  	
  	If $S$ is linearly dependent, then some vector in $S$ can be written as a linear combination of the others. By part (a), we can remove that vector from $S$ without changing $H$. Repeat until we obtain a linearly independent set from $S$. Since $H$ is non-zero, this process must end with some number of vectors from $S$. 	
  \end{proof}
  
  \Answer
  
  \begin{proof}
  	Let $U$ be the reduced row echelon form of $A$. The set of pivot columns of $U$ is linearly independent since no column is a linear combination of the columns that proceed it (by first Theorem today).
  	
  	Since $A$ is row equivalent to $U$, the pivot columns of $A$ are also linearly independent.
  	
  	Recall $\text{Col}(A)=\text{Span}\{a_1, \dots, a_n\}.$ If $a_k$ is a non-pivot column, then it is a linear combination of the other columns. By the spanning set theorem, we can remove each non-pivot column. Thus the pivot columns form a basis for the column space.
  \end{proof}
  
  \Answer The first two columns are pivot columns, so by the previous theorem, we get $$\text{Col}(A)=\text{Span}\left\{\begin{bmatrix} 1\\ 0\\ 0\\ \end{bmatrix}, \begin{bmatrix} 0\\ 2\\ 0\\ \end{bmatrix}  \right\}$$ and a basis is given by $$\left\{\begin{bmatrix} 1\\ 0\\ 0\\ \end{bmatrix}, \begin{bmatrix} 0\\ 2\\ 0\\ \end{bmatrix}  \right\}.$$
  \Answer No, the vectors are not linearly independent. Note: $(1+t)+(1-t)=2$ is a dependence relation. On the other hand, $(1+t)$ and $(1-t)$ do form a basis for $\mathbb{P}_1$. We would just need to check that they span $\mathbb{P}_1$ and they are linearly independent.

\end{enumerate}  

\end{comment}
%\end{document}


%%% Local ables: 
%%% mode: latex
%%% TeX-master: t
%%% End: 
